\documentclass[10pt,a4paper]{amsart}
\usepackage[margin=19mm]{geometry}
\usepackage{amsmath,amssymb,mathtools}
\usepackage[scr=rsfs,cal=euler]{mathalfa}
\usepackage{xfrac}
\usepackage[unicode,hidelinks,bookmarks=false]{hyperref}
\newcommand{\ie}{i.e.\ }
\newcommand{\cf}{cf.\ }
\newcommand{\resp}{respectively\ }
\newcommand{\Subsection}{\subsection}
\providecommand{\nobreakdash}{\nobreak\hbox{-}}
\setlength{\emergencystretch}{2em}
\multlinegap0pt
\usepackage{amssymb}
\usepackage{mathtools}
\usepackage{xcolor}
\usepackage{bm}
\usepackage[shortlabels]{enumitem}
\usepackage[normalem]{ulem}
\newtheorem{lemma}{Lemma}
\newtheorem{claim}{Claim}
\newtheorem{claimproof}{Claimproof}
\newcommand{\BU}{\Bb} %Basis uniformity
\newcommand{\Bb}{\mathcal{B}}
\newcommand{\Ff}{\mathcal{F}}
\newcommand{\Mm}{\mathcal{M}}
\newcommand{\X}[1]{X^{(#1)}}
\title{NTP2 topological structures}
\author{Pablo And\'ujar Guerrero}
\date{Source: arXiv:2604.24522v1 (CC BY 4.0)}
\begin{document}
\maketitle

\noindent\emph{Source and licence.} NTP2 topological structures. Version arXiv:2604.24522v1 (CC BY 4.0), distributed by arXiv under the Creative Commons Attribution 4.0 International licence, http://creativecommons.org/licenses/by/4.0/, as recorded in the arXiv licence metadata for this exact version and shown on the abstract page https://arxiv.org/abs/2604.24522v1. Full licence notice: \texttt{licenses/arxiv-2604.24522-CC-BY-4.0.txt}.\par\smallskip

\noindent\textit{Editorial scope.} Counted unit: the article's lemma lem:closed-chain (source0.tex lines 891--898). The article's complete original proof of it is delivered with it (source0.tex lines 899--919, one \texttt{proof} environment of 21 lines). No mathematics is written, repaired or re-derived: the statement and the proof are the article's own lines, in the article's own notation, and no retained line is substituted or re-flowed.  Retained uncounted as the source-local context the proof needs: lines 217--217.  Nothing was omitted from any retained span, so the package carries no omission record.  No retained line contains a citation, so the wrapper carries no bibliography and no bibliography entry.  No retained line contains a figure environment, an image-inclusion command or drawing code, so no image is delivered and the package declares no asset dependency.  Editorial formatting only, in addition: an AMS \texttt{amsart} preamble that restates the article's own declarations and macros used by the retained lines, namely the environments \texttt{lemma}, \texttt{claim}, \texttt{claimproof} on separate counters, together with the standard packages those lines need.  The retained environments print in this document as Lemma 1, Claim 1, Claimproof 1 in order of appearance, whereas the article prints the same items with the numbering of its own full document.  The complete native source of the article as distributed in the arXiv e-print for this exact version is bundled unchanged as \texttt{sources/199232/source0.tex}.
\section{Preliminaries} \label{sec:prelim}

\begin{lemma}\label{lem:closed-chain}
Let $\Mm$ be a uniform structure. Let $\kappa>1$ be a cardinal and $\{ \X{\alpha} : \alpha < \kappa\}$ be a family of nonempty definable subsets of $M^d$ satisfying that, for all $\alpha < \beta < \kappa$,
\begin{enumerate}
    \item $\X{\beta} \subseteq \X{\alpha}$,
    \item $\X{\beta}$ is nowhere dense in $\X{\alpha}$ (in the subspace topology).
\end{enumerate}
Then $\X{0}$ admits an inp-pattern of depth $\kappa-1$. 
\end{lemma}

\begin{proof}
Let $\BU$ be a definable basis for the product uniformity on $M^d$. Let $\{ \X{\alpha} : \alpha < \kappa\}$ be as in the lemma. By passing to the closures if necessary, we may assume that every set $\X{\alpha}$, for $0<\alpha<\kappa$, is closed in $\X{0}$. For each $0<\alpha<\kappa$, let $\varphi_\alpha$ be a (partitioned) formula that defines the family of sets of the form $(U\setminus V)[\X{\alpha}]$, for $U,V\in \BU$. We prove the existence of an inp-pattern of depth $\kappa-1$ witnessed by the formulas $\varphi_\alpha$. Specifically, we describe, for every $m>1$ and $0<\alpha_1<\cdots\alpha_m < \kappa$, entourages $U_{i,j}$, for $1\leq i\leq m$ and $1\leq j\leq m+1$, such that the family of sets $\{ (U_{i,j} \setminus U_{i,j+1})[\X{\alpha_i}] : i, j\leq m\}$ forms an inp-pattern of depth $m$ and length $m$ in $\X{0}$ (namely conditions (1) and (2) in the definition of inp-pattern in Section~\ref{sec:prelim} are satisfied, and moreover the families described in condition (2) are consistent with $\X{0}$). For simplicity of notation, and since there will be no room for confusion, we set $\alpha_i=i$ for every $i\leq m$.

Let us fix $m>1$. For any $n\leq m$, suppose that we have defined entourages $U_{i,j}$ for $n\leq i \leq m$ and $j\leq m+1$, and consider the following two statements.
\begin{description}
    \item[$(i)_n$] Each family $\{ (U_{i,j} \setminus U_{i,j+1})[\X{i}] : j\leq m\}$, for $n\leq i\leq m$, is pairwise disjoint. 
    \item[$(ii)_n$] For every map $\eta:\{n,\ldots, m\}\rightarrow \{1,\ldots, m\}$, the intersection of the family $\{ (U_{i,\eta(i)} \setminus U_{i,\eta(i)+1})[\X{i}] : n\leq i\leq m \}$ has nonempty interior which intersects $\X{n-1}$.
\end{description}
We prove the lemma by proving $(i)_1$ and $(ii)_1$. We proceed by reverse induction on $n\leq m$. We will use a simple claim. 

\begin{claim}\label{claim:UV}
Let $\Ff$ be a finite family of nonempty open sets in $M^d$, each intersecting $\X{n}$. For any entourage $U\in \BU$, there exists another entourage $V\in \BU$ such that, for all $W\in \Ff$, the \textbf{interior} of the set $W \cap (U\setminus V)[\X{n}]$ intersects the set $\X{n-1}$.
\end{claim}
\begin{claimproof}
Since $\X{n}$ has empty interior in $\X{n-1}$, for each $W\in \Ff$ there exists a point $x_W \in W\cap U[\X{n}] \cap  \X{n-1} \setminus \X{n}$. Since $\X{n}$ is closed, there exists $V \in \BU$ such that $\{x_W : W\in \Ff\}$ is disjoint from the closure of $V[\X{n}]$. The claim follows. 
\end{claimproof}

Let $U_{m,1}\in \BU$ be any entourage. By repeated application of Claim~\ref{claim:UV} (with $\Ff=\{M^d\}$) there exists a decreasing sequence of entourages $U_{m,j}$, for $j\leq m+1$, satisfying that the interior of each set $(U_{m,j} \setminus U_{m,j+1})[\X{m}]$ intersects the set $\X{m-1}$. In particular condition $(ii)_m$ is clearly satisfied. Since the sequence of entourages is decreasing, condition $(i)_m$ is also immediate. 

We now fix $n<m$, and by induction hypothesis we assume that both $(i)_{n+1}$ and $(ii)_{n+1}$ hold, witnessed by entourages $U_{i,j}$, for $n<i\leq m$ and $j\leq m+1$. For each map $\eta:\{n+1,\ldots, m\}\rightarrow \{1,\ldots, m\}$, let us denote the interior of the intersection $\bigcap \{(U_{i,\eta(i)} \setminus U_{i,\eta(i)+1})[\X{i}] : n< i\leq m \}$ by $W_\eta$. By condition $(ii)_{n+1}$, every such open set $W_\eta$ intersects $\X{n}$. Let $U_{n,1}\in \BU$ be any entourage. By repeated application of Claim~\ref{claim:UV}, we may find a decreasing sequence of entourages $U_{n,j}$, for $j\leq m+1$, satisfying that, for any $\eta:\{n+1,\ldots, m\}\rightarrow \{1,\ldots, m\}$ and $j\leq m$, the interior of the set $W_\eta \cap (U_{n,j} \setminus U_{n,j+1})[\X{n}]$ intersects the set $\X{n-1}$. Condition $(ii)_n$ clearly follows. Furthermore, since once again the entourages $U_{n,j}$ are decreasing, condition $(i)_n$ is immediate. 
\end{proof}

\end{document}
